\documentclass[UTF8,fontset=fandol,11pt,a4paper]{ctexart}
\usepackage[margin=22mm,headheight=15pt]{geometry}
\usepackage{xcolor,booktabs,longtable,tabularx,array,enumitem,fvextra,tcolorbox,fancyhdr}
\usepackage[unicode,colorlinks=true,linkcolor=blue!45!black,urlcolor=blue!55!black]{hyperref}
\usepackage{xurl}
\hypersetup{pdftitle={NetworkGuardian 使用说明},pdfauthor={NetworkGuardian 项目},pdfsubject={Windows 网络管理与本机 HTTP MCP}}
\setlength{\parindent}{0pt}
\setlength{\parskip}{5pt}
\setlist{itemsep=3pt,topsep=5pt}
\pagestyle{fancy}
\fancyhf{}
\fancyhead[L]{\small NetworkGuardian 使用说明}
\fancyhead[R]{\small 0.13.0}
\fancyfoot[C]{\thepage}
\newcommand{\ui}[1]{\textbf{“#1”}}
\newcommand{\tool}[1]{\nolinkurl{#1}}
\newtcolorbox{guidebox}[1]{colback=blue!3,colframe=blue!35!black,title=#1,fonttitle=\bfseries,boxrule=0.5pt}
\DefineVerbatimEnvironment{Code}{Verbatim}{fontsize=\small,breaklines=true,breakanywhere=true,samepage=true,frame=single,framesep=6pt}
\begin{document}
\begin{titlepage}
\centering
\vspace*{22mm}
{\Huge\bfseries NetworkGuardian\par}
\vspace{8mm}
{\LARGE 使用说明\par}
\vspace{5mm}
{\large Windows 单文件版 · 0.13.0\par}
\vspace{17mm}
\begin{minipage}{0.82\textwidth}
这份说明从第一次打开软件开始，介绍如何查看网络是否正常、连接无线网络、指定网卡、保持禁用设备、理解本地模型，以及让 AI 通过本机 MCP 服务器使用软件。

不需要先懂路由、无线驱动或机器学习。高级参数与工具名称放在后半部分；首次使用可先阅读第 1--4 节。
\end{minipage}
\vfill
{\large 2026 年 10 月 16 日\par}
\vspace{5mm}
{\small 文档与对应 LaTeX 源码随项目维护；界面名称以本版本为准。\par}
\end{titlepage}
\pagenumbering{Roman}
\begingroup\small\setlength{\parskip}{0pt}\tableofcontents\endgroup
\clearpage
\pagenumbering{arabic}

\section{先了解它能做什么}

NetworkGuardian 是运行在本机的网络管家。它分别查看网卡、无线连接、IP 地址、路由和外网探测结果，在自动恢复开启时按策略尝试恢复连接。
它可以管理多块无线网卡，同时保留有线网络作为出口；也可以记住某个 Wi-Fi 只允许使用哪块网卡。

\begin{guidebox}{三个容易混淆的概念}
\textbf{网卡}是电脑连接网络的设备，例如内置无线网卡、USB 无线网卡或有线网卡。
\textbf{SSID}就是 Wi-Fi 名称。
\textbf{外网出口}是 Windows 当前优先用来访问外网的接口。某块备用网卡能上网，不代表当前出口已经正常。
\end{guidebox}

软件不会把信号强就当作网络好。信号强度只是无线接收情况；真正的质量还取决于丢包、连接失败、延迟和波动。
已经正常工作的连接通常会被保留，避免为追求分数频繁断开重连。

\subsection{准备条件}
\begin{itemize}
\item 使用支持本项目的 Windows 10/11 x64 环境，系统能识别对应的网卡驱动。
\item 主程序是单个 \tool{NetworkGuardian.exe}，不需要另行安装 .NET 运行时。
\item 主程序一般不要求管理员权限；启用、禁用、重启设备或运行管理员命令时，Windows 可能显示 UAC 提权提示。
\item Npcap 用于更准确地验证探测实际经过哪块网卡。未安装时软件会提示并降级；安装入口为 \url{https://npcap.com/}。
\end{itemize}
软件不会替你提供校园网账号、路由器密码或驱动。需要认证的网络，仍需配置自己的有效凭据或认证程序。

\section{第一次启动：按这个顺序操作}
\begin{enumerate}
\item 将 EXE 放在准备长期使用的位置，双击运行。不要同时启动新旧两个版本。
\item 如果没有看到窗口，查看任务栏右下角的托盘区域；也可能藏在向上的小箭头内。点击 NetworkGuardian 图标打开窗口。
\item 先看\ui{运行概览}顶部的\ui{网络质量评价报告}。第一次使用没有足够样本，显示“等待采样”是正常情况。
\item 打开\ui{网卡管理}，核对设备名称。不要的网卡可以设置\ui{保持禁用}；需要使用的网卡应保持启用。
\item 打开\ui{无线网络}，点击扫描，查看各网卡能看到哪些 Wi-Fi。没有凭据时，先到\ui{网络凭据}添加网络。
\item 确认网络选择和校园网信息正确后，到\ui{自动管理}检查自动恢复、自动启用和启动方式。
\end{enumerate}

设置页的改动会在短暂防抖后自动保存；\ui{保存并应用}可以立即确认保存结果。不要把页面切换当成“取消未保存改动”。

\subsection{侧栏怎么分}
\begin{tabularx}{\linewidth}{@{}p{23mm}p{29mm}X@{}}
\toprule
分类 & 入口 & 适合在这里做什么\\
\midrule
总览 & 运行概览 & 阅读质量报告、查看出口与任务、回答偏好选择题。\\
网络管理 & 无线网络 & 扫描、连接或断开 Wi-Fi，设置校园归属与指定网卡。\\
网络管理 & 有线网络 & 查看有线接口、IP、路由和外网探测。\\
网络管理 & 网卡管理 & 启用设备、保持禁用、取消离线设备的禁用规则。\\
网络管理 & 网络凭据 & 保存自己的 Wi-Fi 账号、密码及高级认证配置。\\
策略与诊断 & 自动管理 & 设置恢复、探测、质量学习、认证程序、启动和 MCP。\\
策略与诊断 & 诊断日志 & 按等级、分类和关键词查询发生过什么。\\
\bottomrule
\end{tabularx}

\subsection{哪些操作会再次确认}
切换页面、查看记录、筛选日志、刷新和扫描无需确认。
连接、断开、写入设置、提交回答、删除、导出和启用/禁用设备等界面按钮会先确认，默认选择取消。
\textbf{MCP 写工具按 AI 客户端发出的有效调用直接执行，不弹出本软件的界面确认。}Windows 自己的 UAC 提权仍可能出现。

关闭窗口可以隐藏到托盘，也可以退出，取决于启动设置。隐藏到托盘后，后台监测和已开启的 MCP 服务仍在运行；从托盘选择退出才会结束进程。

\section{阅读网络质量评价报告}
报告会自动随状态刷新，按以下顺序阅读即可：
\begin{enumerate}
\item \textbf{是否能上网。}“目前可以上网”是近期外网探测的结论；“暂时还不能判断”表示没有新结果，不等于已经断网。
\item \textbf{当前出口。}核对现在走有线、无线还是其他接口。备用网络可用但当前出口失败时，报告会单独指出。
\item \textbf{各无线连接的评价。}看具体设备和 Wi-Fi 名称，避免把一块网卡的历史成绩当成另一块网卡的质量。
\item \textbf{建议。}例如保留当前可用出口、继续观察、排查本地链路，或核对认证和路由。
\end{enumerate}

\begin{guidebox}{一段报告应当怎么理解}
“某无线连接本地链路有丢包，可能出现卡顿；建议保留当前可用出口，再检查信号、干扰或网卡链路。”

这说明探测发现了本地链路风险，\textbf{不等于}已经证明 USB 网卡损坏，也不表示软件马上要切断当前网络。
\end{guidebox}

\begin{tabularx}{\linewidth}{@{}p{27mm}X@{}}
\toprule
用词 & 简单解释\\
\midrule
丢包 & 发出的探测包没有全部收到回应；需要结合协议和回应历史判断。\\
延迟 & 得到回应等待多久，通常用毫秒表示；实时通话和游戏更在意它。\\
抖动 & 延迟忽高忽低的程度；平均延迟不高也可能因波动而卡顿。\\
TCP 失败 & 外网连接尝试失败的比例，比单独看公网 ping 更有参考价值。\\
依据充分度 & 当前有效样本有多少，不能解释成操作成功概率。\\
质量分数 & 软件按近期探测计算的综合参考，不是下载速度，也不是网卡硬件评分证书。\\
\bottomrule
\end{tabularx}

至少五组、五分钟内且与当前网卡、SSID、频段一致，并有有效本地或外网连接测量的数据，才会给出评级。
样本不足、指标未知、未连接、历史过期或关闭采样时，报告会说明限制。
部分设备不回复 ping；如果 TCP 连接仍正常，不会仅因此判定断网。
当前没有持续的有线质量评分和下载带宽测量，报告不会据此虚构有线分数或下载速度。

\section{连接、绑定和禁用网卡}
\subsection{连接一个 Wi-Fi}
\begin{enumerate}
\item 到\ui{无线网络}扫描，确认是哪块网卡看到目标 Wi-Fi。
\item 如果已有系统配置，可以连接；如果需要新密码或校园网账号，到\ui{网络凭据}添加。
\item 连接后回到概览，看接口是否已连接、是否有地址、外网探测是否通过。
\end{enumerate}
“已连接 Wi-Fi”“拿到 IP”和“外网可用”是三个不同阶段。连接成功但还不能上网时，先看认证与探测结果，不要立即认定网卡坏了。

\subsection{让某个 Wi-Fi 只用指定网卡}
在无线网络列表中点击该 SSID 的\ui{指定网卡}，从设备名称和接口别名中选择目标。
这对校园网和普通 Wi-Fi 都有效；SSID 指定优先于质量评级和学习偏好。
选择\ui{不指定（自动评级分配）}即可取消。被指定的设备如果离线或保持禁用，软件不会为了连接该 SSID 擅自改用其他网卡。

\subsection{保持禁用：不让自动启用把它打开}
\begin{enumerate}
\item 打开\ui{网卡管理}，核对设备名称与类型。
\item 点击\ui{保持禁用}，确认断开该设备连接的影响；需要时处理 Windows 的提权提示。
\item 界面会分别说明规则是否保存、设备操作是否完成。如果失败，规则仍保留，之后会退避重试。
\end{enumerate}
软件运行期间，保持禁用优先于自动启用、故障重启、自动选网及软件无线电开启。
即使暂停自动恢复，已保存的保持禁用仍会执行；设备被其他程序重新启用后，会在设备通知或巡检中再次尝试禁用，不是内核级实时拦截。

需要恢复时点击\ui{取消禁用并启用}。设备拔下后，可在离线规则中取消禁用。
规则按设备标识匹配；更换端口或驱动导致设备标识变化时，可能需要重新指定。
如果 Windows 要求重启，界面会明确提示，不把“需要重启”当成已经完成。

\subsection{断开、暂停和禁用的区别}
\textbf{断开}只结束一次 Wi-Fi 连接，自动恢复可能随后重连。
\textbf{暂停自动恢复}暂停自动恢复动作，不等于禁用设备。
\textbf{保持禁用}保存长期设备规则，并持续维持禁用。
如果只是想暂时观察，不要误用长期禁用；如果明确不用某块网卡，也不要只反复点断开。

\section{校园网与网络凭据}
\subsection{两种常见校园网}
\textbf{需要浏览器或独立客户端认证的网络：}先连接 Wi-Fi 或网线，再完成学校要求的认证流程。需要软件调用认证程序时，到自动管理配置程序路径、参数及触发条件。

\textbf{使用 802.1X 账号认证的 Wi-Fi：}在网络凭据中选择企业网络，填写学校提供的身份、密码、EAP 方法及证书要求。不要把它当成普通家用 Wi-Fi 密码。
PEAP、EAP-TLS 和自定义 XML 面向不同部署；不清楚时先向学校取得正确参数，不要凭猜测关闭证书校验。

\subsection{保存与应用}
凭据库独立于 Windows 的 Wi-Fi 配置。企业网络通常需要向对应接口写入系统配置和用户认证数据；个人与开放网络可以使用临时 WLAN 连接方式。
\ui{写入到所有无线网卡}表示准备配置，不表示已经连接或通过外网探测。
删除凭据库条目不会自动删除已经写入 Windows 的配置；需要移除系统配置时使用对应删除操作。

保存密码使用当前 Windows 用户的 DPAPI 保护。复制数据到另一台电脑或另一个用户后，密码可能无法解密，需要重新输入；单独复制 EXE 也不会带走原来的数据目录。

\subsection{使用外部认证命令}
在自动管理中配置你已有的认证程序或离线命令。区分程序路径和参数，只有确实需要时才启用管理员运行。
超时、冷却和每小时次数用于避免反复弹窗。通过 MCP 的 \tool{command_run} 可显式运行一个已启用的配置项；校园认证固定使用 \tool{campus-auth} 作为 ID。
工具返回“已启动”不等于认证成功；等待后续外网验证再判断。

\section{本地模型如何学习}
当前模型是一个 8--8--1 的小型数值神经网络，输入主要是质量、失败率、延迟、波动、信号、证据充分度和联网情况。
它可以在已有规则内有限调整选网偏好，并从有实测支持的回答中继续训练。它不是能理解任意文本并自由对话的大语言模型。

概览偏好卡提供约 60 道动态选择题，题目会结合当前 SSID 和设备生成。可以回答，也可以直接换题；跳过不会当成训练回答。
回答时记录出题及作答时的网络状态，并结合探测降低不可靠回答的权重。
不是每题都会增加神经网络训练次数：校园归属、成本和角色等先作为软标签；有足够实测的网卡偏好或质量取舍才可能直接更新数值模型。

模型不能绕过 SSID 指定、保持禁用或实测可用性约束。训练次数增加不等于准确率已提高，也不表示会立即换网。
质量报告的中文由实测解释规则生成，不是神经网络自由编造诊断；运行期不依赖在线教师模型或 API key。

需要管理个人数据时，使用查看记录、删除、导出或重置学习。删除会按保留数据重训；重置恢复内置模型。
答题集通常最多保留 1,000 条，不能把长期运行理解成无限保存全部行为。

\section{让 AI 通过本机 MCP 使用软件}
\subsection{开启服务}
\begin{enumerate}
\item 打开\ui{自动管理}中的\ui{AI 工具接入 / MCP}。
\item 打开启用开关。默认关闭；默认固定端口为 \tool{18765}，也可以设置 1--65535 范围内的其他端口。
\item 等待显示\ui{正在监听}。默认完整地址是：
\begin{Code}
http://127.0.0.1:18765/mcp
\end{Code}
\item 在本机 AI 客户端添加一个支持 \textbf{Streamable HTTP} 的 MCP 服务器，填入完整地址。
\item 先让 AI 查询状态和设备列表，再进行控制。
\end{enumerate}
端口被占用时会显示错误，不会偷偷换端口或监听其他网卡。改端口后，要同步修改客户端地址。
仅隐藏到托盘不会关闭服务；停用开关或退出程序会停止监听。

\begin{guidebox}{连接范围与执行方式}
服务器只监听 IPv4 地址 \tool{127.0.0.1}，没有账号、密码、令牌或鉴权。
同一台电脑上能够访问该地址的程序可以调用工具；有效的写工具会直接执行。
不能设置成局域网监听地址。云端服务直接访问自己的 \tool{127.0.0.1}，并不能访问你的电脑；需要由本机 MCP 客户端建立连接。
\end{guidebox}

不同 AI 客户端配置格式不完全相同。下列 JSON 仅表示常见字段结构；以你的客户端说明为准，不要额外配置 stdio 启动命令或旧式 SSE 地址：
\par\noindent\begin{minipage}{\linewidth}
\begin{Code}
{
  "mcpServers": {
    "NetworkGuardian": {
      "type": "http",
      "url": "http://127.0.0.1:18765/mcp"
    }
  }
}
\end{Code}
\end{minipage}

\subsection{第一次可以这样问 AI}
\begin{enumerate}
\item “查询 NetworkGuardian 当前状态和网络质量报告，用简单中文解释，不修改设置。”
\item “列出物理无线网卡和各自连接的 Wi-Fi，告诉我哪条连接存在丢包证据。”
\item “读取当前偏好选择题和选项；我给出答案后再提交，不要替我猜校园网归属。”
\item “把我指定的 Wi-Fi 绑定到查询结果中的这块网卡，然后重新读取配置核对。”
\end{enumerate}
工具需要的 GUID、设备 ID、凭据 ID 和题目 ID 应来自查询结果，不能根据显示名称猜测。
设备状态、题目和可选网络会变化，旧 ID 或旧题目失效时应重新查询。

\subsection{配置写入规则}
\tool{config_get} 可返回全部设置，也可用 \tool{section} 读取顶层分组，如 \tool{wifi}、\tool{quality}、\tool{campusAuth} 或 \tool{mcp}。
\tool{config_patch} 递归合并对象，\textbf{数组整体替换}；未知字段、错误类型和非法值会拒绝，\tool{version} 为只读字段。
例如只调整质量切换收益阈值，不覆盖其他设置：
\begin{Code}
{
  "patch": {
    "quality": { "minimumGain": 24 }
  }
}
\end{Code}
使用实际\ui{读取结果}核对字段名和当前值。修改命令数组前先完整读取，否则整体替换会删除未包含的旧命令。
通过 MCP 修改 MCP 自身端口或开关时，本次结果仍返回；下一次调用要使用新地址，关闭后则不能继续调用。

\subsection{模型投喂不是上传随意文字}
\tool{learning_question} 返回真实题目和选项；\tool{learning_answer} 提交 \tool{questionId} 与选项 \tool{value}。
软件捕获当前网络快照，并按现有可信度规则记录和学习。
\tool{learning_options} 与 \tool{learning_choose} 用于当前网络偏好选择。
这些接口不会把 AI 编造的丢包率当成实测数据，也没有“上传一段文本就训练通用大模型”的功能。

\subsection{理解工具返回值}
\begin{itemize}
\item MCP 的 \tool{isError=true} 表示工具操作失败；阅读文字或结构化结果中的原因。
\item 设备返回 \tool{ruleSaved=true} 但操作失败，表示规则已保存，设备状态还未达到目标。
\item \tool{accepted=true, completed=false} 表示只受理了刷新，不表示新状态已经到达。
\item 命令启动、配置写入和 Wi-Fi 关联成功都不能替代外网探测结果。
\end{itemize}


\subsection{开发者的最小连接示例}
只使用 AI 客户端的人无需执行下面的代码。开发者安装 Python 的官方 \tool{mcp} 包后，可以查询服务；这段代码不会修改网络：
\begin{Code}
import asyncio
from mcp import ClientSession
from mcp.client.streamable_http import streamablehttp_client

async def main():
    async with streamablehttp_client(
        "http://127.0.0.1:18765/mcp",
        terminate_on_close=False,
    ) as (read, write, _):
        async with ClientSession(read, write) as session:
            await session.initialize()
            result = await session.call_tool("quality_report", {})
            print(result.structuredContent)

asyncio.run(main())
\end{Code}
本实现使用 JSON 响应的 Streamable HTTP，支持 2025-11-25、2025-06-18 和 2025-03-26 协议版本；不提供旧式独立 SSE 地址。
\tool{GET /mcp} 返回 405 是本实现不提供服务端 SSE 流的正常结果，不是“服务没有启动”。
\tool{GET /health} 可查询监听状态。MCP 连接仍应填写完整的 \tool{/mcp} 地址。

\section{常见问题与调试顺序}
\subsection{界面没有看到网卡，或扫描为空}
先在网卡管理检查是否保持禁用、驱动是否工作，再刷新设备。
无线扫描还可能受 Windows 位置权限、无线电开关及驱动限制影响；按界面提示处理权限，不要把“扫描被拒绝”当成没有附近 Wi-Fi。
虚拟网卡默认过滤，可在设备页切换显示；虚拟接口不会被当作物理设备启用或禁用。

\subsection{连接了 Wi-Fi，却仍不能上网}
检查是否取得地址、是否需要校园认证、当前出口走哪块网卡，再运行连通性测试。
使用 VPN、TUN、代理或 DNS 接管时，查看概览兼容说明。
某些 VPN 的 kill switch 会禁止物理网卡直连；此时物理网卡探测失败和代理内网页可用可能同时发生，不能只凭一项结果归因。

\subsection{设置禁用后没有立即禁用}
先看设备操作是否失败、是否等待 UAC、是否要求重启。
规则已保存不等于驱动已完成禁用；失败时会保留规则并延迟重试。
要恢复使用，取消保持禁用，不要只去 Windows 临时启用后又与管家策略反复竞争。

\subsection{模型更新少，或者网络没有变化}
回答题目与模型参数更新不是一一对应；样本类型、有效测量和可信度都影响训练。
稳定连接通常不会为了小幅评分优势而切换。先查看问题记录、个人更新次数和质量证据，不要用“是否立刻换网”判断模型是否生效。

\subsection{AI 连不上 MCP}
按顺序检查：软件仍在运行；MCP 开关开启；状态显示正在监听；客户端地址、端口和 \tool{/mcp} 路径正确；客户端支持 Streamable HTTP；客户端运行在本机。
如显示端口占用，选择空闲端口，并更新客户端。端口变更后旧连接不能继续使用。
来自非本地 Origin 或错误 Host 的请求会被拒绝；不要通过远程反向代理把它当作局域网服务使用。

\subsection{让 AI 调试，但先不动网络}
可以要求 AI 依次调用 \tool{guardian_status}、\tool{quality_report}、\tool{debug_bundle}、\tool{debug_logs} 和 \tool{debug_decision}，根据时间、接口、失败原因解释现状。
\tool{debug_model} 用于确认模型结构与训练状态，不是修改权重的入口。
调试查询失败时，应记录工具名、错误和时间，不要让 AI 无限重试网卡重启或认证命令。

\section{数据、更新和备份}
默认数据目录为：
\begin{Code}
%LOCALAPPDATA%\NetworkGuardian\
\end{Code}
\begin{tabularx}{\linewidth}{@{}p{58mm}X@{}}
\toprule
文件或目录 & 内容\\\midrule
\tool{config.json} & 设置、SSID 指定、保持禁用及 MCP 开关与端口。\\
\tool{wifi-networks.json} & 无线凭据库；密码使用当前用户保护。\\
\tool{network-history.json} & 质量历史、偏好参数和权威模型状态。\\
\tool{network-brain.json} & 可查看的模型权重副本。\\
\tool{network-survey.json} & 答题集与上下文。\\
\tool{Logs} & 按设置滚动写入的文件日志。\\
\bottomrule
\end{tabularx}

更新时，先从托盘退出旧版，再替换或运行新版 EXE。不要仅关闭主窗口后认为旧版已退出。
同一用户和默认数据目录下，原来的设置及学习数据会继续使用。
备份前最好退出程序，再复制整个数据目录；不要只备份权重副本而遗漏历史和问答集。
导出的问答、日志和配置可能含 SSID、设备名称或用户自己填写的命令参数，交给他人排查前按实际需要检查内容。

\section{文档与开发验证说明}
本说明对应 0.13.0 的功能。项目提供单元测试和隔离配置的 MCP 检查脚本，后者通过官方 Python MCP 客户端连接真实单文件成品。
隔离检查验证查询、配置、答题与凭据流程，不为了测试而禁用用户正在使用的真实网卡。
硬件驱动、校园认证服务和 UAC 的实际结果仍以当前电脑的返回与后续探测为准。

说明书源码为本文件，PDF 随源码提交。协议参考：
\begin{itemize}
\item MCP Streamable HTTP：\url{https://modelcontextprotocol.io/specification/2025-11-25/basic/transports}
\item MCP 工具：\url{https://modelcontextprotocol.io/specification/2025-11-25/server/tools}
\end{itemize}
\clearpage
\appendix
\section{MCP 工具索引}
下面的“写”表示可能修改配置、连接、学习数据或设备；“查”包括主动扫描和连通性探测。
所有工具的参数以 \tool{tools/list} 返回的 JSON Schema 为准。

\begingroup\footnotesize\setlength{\parskip}{0pt}
\begin{longtable}{@{}p{48mm}p{9mm}p{96mm}@{}}
\toprule 工具名 & 类型 & 用途\\\midrule\endfirsthead
\toprule 工具名 & 类型 & 用途\\\midrule\endhead
\tool{guardian_status} & 查 & 运行快照、接口、路由与服务器状态。\\
\tool{quality_report} & 查 & 中文评价及质量实测数据。\\
\tool{devices_list} & 查 & 设备标识、分类和保持禁用列表。\\
\tool{device_set_state} & 写 & \tool{enabled} 启用；\tool{kept_disabled} 保持禁用；\tool{release} 只解除规则。\\
\tool{device_restart} & 写 & 重启物理设备修复驱动；保持禁用时拒绝。\\
\tool{wifi_list} & 查 & 无线接口、连接、配置及最近扫描。\\
\tool{wifi_scan} & 查 & 扫描指定接口或全部接口，返回完成后的结果。\\
\tool{wifi_connect} & 写 & 按接口连接已保存配置或凭据库网络。\\
\tool{wifi_disconnect} & 写 & 断开无线连接，可能随后被自动恢复。\\
\tool{wifi_radio_set} & 写 & 改变软件无线电开关；看门狗可能重新打开。\\
\tool{wifi_bind} & 写 & 设置 SSID 指定；省略 adapter 取消。\\
\tool{campus_set} & 写 & 设置或取消 SSID 校园网标记。\\
\tool{config_get} & 查 & 全部配置或指定顶层分组。\\
\tool{config_patch} & 写 & 校验后合并、保存并应用设置。\\
\tool{config_reload} & 写 & 重新加载磁盘配置。\\
\tool{guardian_pause} & 写 & 暂停或恢复自动恢复。\\
\tool{guardian_refresh} & 查 & 请求立即刷新枚举与探测。\\
\tool{connectivity_test} & 查 & 立即执行外网探测，可能需要数十秒。\\
\tool{recovery_reset} & 写 & 清除失败降权与冷却。\\
\tool{command_run} & 写 & 按已有配置 ID 执行认证或离线命令。\\
\tool{credentials_list} & 查 & 默认隐藏密码；显式设置 includeSecrets 可读取内存密码。\\
\tool{credentials_upsert} & 写 & 新增或更新凭据；password 单独传入，省略则保留旧值。\\
\tool{credentials_delete} & 写 & 删除库条目，不自动删 Windows 配置。\\
\tool{credentials_apply} & 写 & 将指定库条目应用至网卡；不保证联网。\\
\tool{wifi_profile_delete} & 写 & 删除指定接口或全部接口的系统配置。\\
\tool{learning_status} & 查 & 训练次数、问答记录、偏好和当前题目。\\
\tool{learning_question} & 查 & 获取题目；next=true 跳过并换题。\\
\tool{learning_answer} & 写 & 按真实题目提交答案和现场快照。\\
\tool{learning_options} & 查 & 当前偏好候选及其 key。\\
\tool{learning_choose} & 写 & 选择临时或记忆偏好，可能连接网络。\\
\tool{learning_delete_answer} & 写 & 删除一条回答并重新训练。\\
\tool{learning_reset} & 写 & 重置个人数据及模型。\\
\tool{learning_export} & 写 & 导出本地问答集及偏好数据。\\
\tool{debug_logs} & 查 & 查询最近内存日志，最多 500 条。\\
\tool{debug_clear_logs} & 写 & 清空内存日志，不删除文件日志。\\
\tool{debug_bundle} & 查 & 版本、资源、探测、路由、设备和模型状态汇总。\\
\tool{debug_model} & 查 & 网络权重、有效性与训练计数。\\
\tool{debug_decision} & 查 & 最近实际决策和任务，不重新驱动决策引擎。\\
\bottomrule
\end{longtable}\endgroup
\end{document}
